\documentclass[11pt]{article}
\usepackage{restaurant-docs}

\tikzset{
  entity/.style={
    draw=navy, line width=0.6pt, fill=white, align=left,
    inner sep=5pt, text width=38mm, font=\sffamily\small
  },
  ename/.style={fill=navy, text=white, font=\sffamily\bfseries\small, text width=38mm, align=left, inner sep=5pt},
  actor/.style={draw=navy, fill=teal!12, rounded corners=1pt, font=\sffamily\bfseries\scriptsize, align=center, text width=22mm, minimum height=11mm},
  seqmsg/.style={-{Latex}, teal, thick},
  seqret/.style={-{Latex}, teal, densely dashed}
}

\begin{document}
\docmeta{Technisch ontwerp}

\section{Inleiding en user stories}
Ik bouw een API voor een restaurant. Reserveringen, de menukaart en de rekening
stonden eerst los van elkaar. In deze applicatie horen ze bij één database.
Een klant ziet alleen de eigen reserveringen. Personeel beheert de kaart, de
tafels, de bestellingen en de facturen. Inloggen loopt via Keycloak. De site
\emph{The Table} zit in dezelfde JAR en is Engelstalig.

\begin{description}[style=nextline,leftmargin=1.6cm]
  \item[US1] Als klant wil ik mij registreren en veilig inloggen, zodat ik mijn
    eigen reserveringen kan beheren zonder toegang tot gegevens van andere klanten.
  \item[US2] Als klant wil ik een beschikbare tafel reserveren, mijn reservering
    wijzigen en annuleren, zodat mijn restaurantbezoek passend en actueel is.
  \item[US3] Als personeelslid wil ik bestellingen met meerdere producten beheren
    en daarvoor een factuur met betaalstatus maken, zodat ik de bediening en
    afrekening centraal kan afhandelen.
  \item[US4] Als personeelslid wil ik producten, categorieën, tafels en
    productafbeeldingen beheren, zodat klanten en medewerkers actuele
    restaurantgegevens gebruiken.
\end{description}

Facturen en tafels zijn ondersteunend aan reserveren en bestellen. Naast accounts
bevat het systeem drie operationele gebieden: reserveringen, bestellingen/facturen
en menukaartbeheer met upload en download.

\section{Architectuur}
De code staat in vier lagen:

\begin{itemize}
  \item \textbf{Controller} vertaalt HTTP naar gevalideerde DTO-input en zet
        statuscodes. Entiteiten verlaten deze laag niet.
  \item \textbf{Service} beheert transacties, bedrijfsregels en mapping via
        \texttt{Views.of()}.
  \item \textbf{Repository} bevat Spring Data JPA, unieke queries en
        \texttt{PESSIMISTIC\_WRITE}-locks.
  \item \textbf{Model} bevat acht JPA-entiteiten met foreign keys, unieke
        constraints en versienummers.
\end{itemize}

\texttt{open-in-view} staat uit. JWT-validatie controleert handtekening, issuer,
audience \texttt{restaurant-api} en subject. H2 is de standaarddatabase;
PostgreSQL is optioneel via omgevingsvariabelen.

\section{Functionele eisen --- klant}
\begin{longtable}{@{}>{\bfseries}p{1.2cm}p{1.6cm}p{11.2cm}@{}}
\toprule
Eis & Story & Acceptatiecriterium \\
\midrule
\endhead
F01 & US1 & Nieuwe accounts registreren via Keycloak en krijgen automatisch CUSTOMER. \\
F02 & US1 & Keycloak authenticeert en geeft een ondertekend access token voor restaurant-api. \\
F03 & US1 & Alleen een Keycloak-beheerder kent STAFF toe; zelfregistratie geeft geen personeelsrechten. \\
F04 & US1 & Een klant kan geen reservering van een ander lezen, wijzigen of annuleren (HTTP 403). \\
F05 & US2 & Een ingelogde gebruiker maakt een reservering met naam, tafel, UTC-begintijd, gasten en opmerkingen. \\
F06 & US2 & Alleen een toekomstige begintijd is geldig; eindtijd is begintijd plus twee uur. \\
F07 & US2 & Niet-inzetbare tafels en groepen groter dan de capaciteit worden geweigerd. \\
F08 & US2 & Overlappende actieve reserveringen voor dezelfde tafel worden geweigerd; aansluitende tijdvakken mogen. \\
F09 & US2 & Eigenaar en personeel wijzigen een actieve reservering; geannuleerde reserveringen blijven vergrendeld. \\
F10 & US2 & DELETE annuleert met behoud van historie en geeft het tijdvak vrij. \\
F11 & US2 & GET toont klanten alleen eigen reserveringen en personeel de volledige lijst. \\
\bottomrule
\end{longtable}

\section{Functionele eisen --- personeel}
\begin{longtable}{@{}>{\bfseries}p{1.2cm}p{1.6cm}p{11.2cm}@{}}
\toprule
Eis & Story & Acceptatiecriterium \\
\midrule
\endhead
F12 & US3 & Personeel beheert een bestelling met tafel en 1 tot 100 productregels. \\
F13 & US3 & Een product mag per bestelling eenmaal voorkomen (hoeveelheid 1--100); onbeschikbare producten worden geweigerd. \\
F14 & US3 & De server berekent het totaal uit vastgelegde eenheidsprijzen maal aantallen (\texttt{BigDecimal}). \\
F15 & US3 & Bestaande bestelregels behouden productnaam en eenheidsprijs bij menukaartwijzigingen. \\
F16 & US3 & Maximaal één factuur per bestelling; het bedrag komt uit de bestelling. \\
F17 & US3 & Een gefactureerde bestelling kan niet worden gewijzigd of verwijderd zolang de factuur bestaat. \\
F18 & US3 & Personeel markeert een factuur als betaald; betaald kan niet terug of verwijderd. \\
F19 & US3 & Een onbetaalde factuur mag worden verwijderd, zodat de bestelling weer bewerkbaar is. \\
F20 & US4 & Personeel beheert categorieën en producten; verwijderen alleen indien ongebruikt. \\
F21 & US4 & Personeel beheert tafels; gebruikte tafels en gereserveerde capaciteit zijn beschermd. \\
F22 & US4 & Klanten en personeel lezen catalogus en tafels; mutaties vereisen STAFF. \\
F23 & US4 & Personeel uploadt per product één PNG/JPEG; klanten en personeel downloaden die. \\
\bottomrule
\end{longtable}

\section{Niet-functionele eisen}
\begin{longtable}{@{}>{\bfseries}p{1.2cm}p{1.8cm}p{11cm}@{}}
\toprule
Eis & Story & Acceptatiecriterium \\
\midrule
\endhead
N01 & US1 & JWT-handtekening, issuer, audience en verloopdatum worden gevalideerd; ontbrekend token geeft 401. \\
N02 & US2 & Reserveringscontrole en opslag gebeuren in één transactie met een schrijflock op de tafel. \\
N03 & US3 & Foreign keys, unieke factuur-orderkoppeling, versienummers en numerieke constraints bewaken integriteit. \\
N04 & US4 & Uploads maximaal 2~MiB en 16 miljoen pixels; paden komen nooit uit de bestandsnaam. \\
N05 & US1--4 & Java~21 / Spring Boot met gescheiden controllers, services, repositories en DTO-records. \\
N06 & US1--4 & Bekende fouten als 400/403/404/409/413/415; interne details blijven verborgen. \\
N07 & US2--3 & Maven verify: minstens twintig unit-tests en integratietests; PricingService en ReservationPolicy 100\% line coverage. \\
N08 & US1--4 & JDBC is extern instelbaar; H2 standaard, PostgreSQL-driver meegeleverd. \\
N09 & US4 & \texttt{data.sql} laadt demonstratiegegevens zonder bestaande rijen te overschrijven. \\
\bottomrule
\end{longtable}

\section{Klassendiagram van het domein}
Figuur~\ref{fig:klassen} toont acht entiteiten met kardinaliteiten. Attributen
sluiten aan bij de Java-klassen; methodes zijn weggelaten. \texttt{id} en
\texttt{version} ondersteunen opslag en gelijktijdige wijzigingen.

\begin{figure}[ht]
\centering
\begin{tikzpicture}[
  font=\sffamily,
  box/.style={
    rectangle split, rectangle split parts=2, draw=navy, line width=0.7pt,
    rectangle split part fill={navy,white}, text=navy, inner sep=4pt,
    text width=39mm, font=\scriptsize
  }
]
\node[box] (cat) {\textcolor{white}{\textbf{Category}}
  \nodepart{two} id, version\\name, description};
\node[box, right=8mm of cat] (prod) {\textcolor{white}{\textbf{Product}}
  \nodepart{two} id, version\\name, price, available\\category};
\node[box, right=8mm of prod] (img) {\textcolor{white}{\textbf{ProductImage}}
  \nodepart{two} id, version\\contentType, content\\product};

\node[box, below=16mm of cat] (ord) {\textcolor{white}{\textbf{RestaurantOrder}}
  \nodepart{two} id, version, createdAt\\total, diningTable, lines};
\node[box, below=16mm of prod] (line) {\textcolor{white}{\textbf{OrderLine}}
  \nodepart{two} id, version, quantity\\unitPrice, productName\\order, product};
\node[box, below=16mm of img] (tab) {\textcolor{white}{\textbf{DiningTable}}
  \nodepart{two} id, version\\tableNumber, seats\\available};

\node[box, below=16mm of ord] (inv) {\textcolor{white}{\textbf{Invoice}}
  \nodepart{two} id, version, createdAt\\total, paid, order};
\node[box, below=16mm of tab] (res) {\textcolor{white}{\textbf{Reservation}}
  \nodepart{two} id, version, customerSubject\\customerName, diningTable\\startsAt, endsAt, guests\\remarks, cancelled};

\draw[teal, thick] (cat) -- node[above,font=\tiny\color{teal}]{1} node[above,pos=0.85,font=\tiny\color{teal}]{0..*} (prod);
\draw[teal, thick] (prod) -- node[above,font=\tiny\color{teal}]{1} node[above,pos=0.85,font=\tiny\color{teal}]{0..1} (img);
\draw[teal, thick] (prod) -- node[right,font=\tiny\color{teal}]{1} node[right,pos=0.85,font=\tiny\color{teal}]{0..*} (line);
\draw[teal, thick] (ord) -- node[above,font=\tiny\color{teal}]{1} node[above,pos=0.85,font=\tiny\color{teal}]{1..*} (line);
\draw[teal, thick] (ord) -- node[left,font=\tiny\color{teal}]{1} node[left,pos=0.85,font=\tiny\color{teal}]{0..1} (inv);
\draw[teal, thick] (tab) -- node[right,font=\tiny\color{teal}]{1} node[right,pos=0.85,font=\tiny\color{teal}]{0..*} (res);
\draw[teal, thick] (tab.west) -- ++(-6mm,0) |- (ord.east)
  node[at start,left,font=\tiny\color{teal}]{1}
  node[near end,above,font=\tiny\color{teal}]{0..*};
\end{tikzpicture}
\caption{Domeinmodel: één-op-één voor factuur en productafbeelding; overige relaties één-op-veel.}
\label{fig:klassen}
\end{figure}

\texttt{OrderLine} koppelt bestelling en product met hoeveelheid en historische
prijs. Een geannuleerde reservering blijft bestaan.

\section{Sequentie 1 --- reservering maken}
Figuur~\ref{fig:seq-res} toont \texttt{POST /api/reservations}. Eigenaarschap
komt uit het JWT-subject. De transactie houdt de tafellock vast tot opslag.

\begin{figure}[ht]
\centering
\begin{tikzpicture}[font=\sffamily\scriptsize, x=22mm, y=-9mm]
  \foreach \i/\n in {0/HTTP client,1/Reservation Controller,2/Reservation Service,3/DiningTable Repo,4/Reservation Policy,5/Reservation Repo} {
    \node[actor] (a\i) at (\i,0) {\n};
    \draw[graytext, densely dashed] (a\i.south) -- ++(0,-12.2);
  }
  \draw[seqmsg] (0,-1.2) -- node[above,font=\tiny]{POST input + token} (1,-1.2);
  \draw[seqmsg] (1,-2.1) -- node[above,font=\tiny]{create(input, subject)} (2,-2.1);
  \draw[seqmsg] (2,-3.0) -- node[above,font=\tiny]{findLockedById} (3,-3.0);
  \draw[seqret] (3,-3.9) -- node[above,font=\tiny]{DiningTable / 404} (2,-3.9);
  \draw[seqmsg] (2,-4.8) -- node[above,font=\tiny]{validate(...)} (4,-4.8);
  \draw[seqret] (4,-5.7) -- node[above,font=\tiny]{endsAt / 400, 409} (2,-5.7);
  \draw[seqmsg] (2,-6.6) -- node[above,font=\tiny]{overlaps(...)} (5,-6.6);
  \draw[seqret] (5,-7.5) -- node[above,font=\tiny]{boolean} (2,-7.5);
  \draw[seqmsg] (2,-8.4) -- node[above,font=\tiny]{save (geen overlap)} (5,-8.4);
  \draw[seqret] (5,-9.3) -- node[above,font=\tiny]{Reservation} (2,-9.3);
  \draw[seqret] (2,-10.2) -- node[above,font=\tiny]{ReservationView} (1,-10.2);
  \draw[seqret] (1,-11.1) -- node[above,font=\tiny]{201 + Location} (0,-11.1);
\end{tikzpicture}
\caption{Aanmaken van een reservering met tafellock en overlapcontrole.}
\label{fig:seq-res}
\end{figure}

Alternatieven: geen token 401; ontbrekende tafel 404; ongeldig tijdstip of
capaciteit 400; tafel niet inzetbaar of overlap 409. Bij overlap wordt
\texttt{save()} niet uitgevoerd. De integratietest bewijst één winnaar bij twee
gelijktijdige boekingen.

\section{Sequentie 2 --- factuur maken}
Figuur~\ref{fig:seq-inv} toont \texttt{POST /api/invoices} (alleen STAFF). Het
bedrag komt uit de bestelling, niet uit de requestbody.

\begin{figure}[ht]
\centering
\begin{tikzpicture}[font=\sffamily\scriptsize, x=28mm, y=-9mm]
  \foreach \i/\n in {0/HTTP personeel,1/Invoice Controller,2/Invoice Service,3/Order Repo,4/Invoice Repo} {
    \node[actor] (b\i) at (\i,0) {\n};
    \draw[graytext, densely dashed] (b\i.south) -- ++(0,-10.4);
  }
  \draw[seqmsg] (0,-1.2) -- node[above,font=\tiny]{POST \{orderId\} + token} (1,-1.2);
  \draw[seqmsg] (1,-2.1) -- node[above,font=\tiny]{create(input)} (2,-2.1);
  \draw[seqmsg] (2,-3.0) -- node[above,font=\tiny]{findLockedById} (3,-3.0);
  \draw[seqret] (3,-3.9) -- node[above,font=\tiny]{RestaurantOrder / 404} (2,-3.9);
  \draw[seqmsg] (2,-4.8) -- node[above,font=\tiny]{existsByOrderId} (4,-4.8);
  \draw[seqret] (4,-5.7) -- node[above,font=\tiny]{boolean} (2,-5.7);
  \draw[seqmsg] (2,-6.6) -- node[above,font=\tiny]{save (nog geen factuur)} (4,-6.6);
  \draw[seqret] (4,-7.5) -- node[above,font=\tiny]{Invoice} (2,-7.5);
  \draw[seqret] (2,-8.4) -- node[above,font=\tiny]{InvoiceView} (1,-8.4);
  \draw[seqret] (1,-9.3) -- node[above,font=\tiny]{201 + Location} (0,-9.3);
\end{tikzpicture}
\caption{Aanmaken van een factuur met orderlock en unieke orderkoppeling.}
\label{fig:seq-inv}
\end{figure}

Alternatieven: geen STAFF 403; bestelling ontbreekt 404; bestaande factuur 409.
De unieke \texttt{order\_id}-constraint is een tweede verdedigingslaag.
\texttt{createdAt} komt uit \texttt{Clock}, \texttt{total} uit de bestelling,
\texttt{paid} start als \texttt{false}.

\section{Technische afspraken, verificatie en bronnen}
Verwijderregels beschermen referenties: een categorie met producten, een gebruikt
product en een gebruikte tafel leveren 409. Tijdwaarden zijn UTC zonder suffix.
Bedragen zijn EUR-eindprijzen zonder btw-uitsplitsing.

Verificatie: \texttt{mvn clean verify} met Java~21 gaf 42 tests, 0 fouten
(1~november 2026). JaCoCo bewaakt 100\% line coverage voor
\texttt{PricingService} en \texttt{ReservationPolicy}. Zie
\texttt{docs/verification.md}.

\subsection{Bronnen}
\begin{enumerate}
  \item Ali Khoshkam, \emph{idee fase backend 2.pdf}, aangeleverd lokaal.
  \item NOVI, \emph{Eindopdracht Backend v4.3.pdf}, inclusief quickscan en rubric.
  \item Spring Boot: \url{https://docs.spring.io/spring-boot/3.5/system-requirements.html}
  \item Spring Security JWT: \url{https://docs.spring.io/spring-security/reference/6.5/servlet/oauth2/resource-server/jwt.html}
  \item Keycloak-import: \url{https://www.keycloak.org/server/importExport}
\end{enumerate}

Ik heb dit ontwerp zelf gemaakt. OpenAI Codex heeft me geholpen bij stukken
tekst en uitleg. De prompts die ik zelf heb gestuurd staan in
\texttt{docs/ai-prompts.md}. De diagrammen horen bij de klassen in de broncode.

Broncode: \url{https://github.com/Alikhoshkam1989/myRestaurant}

\end{document}
